\documentclass[12pt,twocolumn]{article}
\usepackage[margin=2cm,columnsep=0.5cm]{geometry}
\usepackage{amsmath,amssymb}
\usepackage{booktabs}
\usepackage{hyperref}
\usepackage{xcolor}

\definecolor{bctnavy}{RGB}{11,37,69}
\definecolor{bctblue}{RGB}{13,71,161}
\definecolor{bctgreen}{RGB}{27,94,32}

\title{\textbf{\color{bctnavy}BCT Letter 114: Plasmoid Physics}\\
{\large\color{bctblue}Macroscopic OHC Hopf Excitations, Acoustic Cavitation,\\
and the Thunderstorm Generator}}
\author{Michel Robert Cabri\'{e}\\
{\small Independent Researcher, Victoria, Australia}\\
{\small ORCID: 0009-0007-9561-9859}\\
{\small \href{mailto:ZeroFreeParameters@gmail.com}
{ZeroFreeParameters@gmail.com}}}
\date{23 March 2026}

\begin{document}
\twocolumn[{%
\maketitle
\begin{center}\rule{0.9\textwidth}{0.4pt}\end{center}
\begin{center}
\parbox{0.88\textwidth}{\small\textbf{Abstract.}
A plasmoid is a self-contained, self-sustaining plasma
structure consisting of two counter-rotating vortices with a
zero-point at their plane of equilibrium. The BCT Superfluid
Lattice Model identifies plasmoids as macroscopic
Octet-Hopfion Condensate (OHC) Hopf charge excitations with
$H = 2$: the simplest stable non-trivial topological
excitation of the quantum vacuum, made visible at macroscopic
scale through acoustic cavitation. The Thunderstorm Generator
of Australian inventor Malcolm Bendall --- a plasmoid-based
device using cavitation in a water bubbler --- is identified
as an accidental BCT Hopf resonator: its geometry is correct
because the geometry of the vacuum demanded it, not because
the inventor derived it from theory. Ball lightning (Letter~111)
is the $H \geq 2$ atmospheric case. The SAFIRE project's
plasma discharge experiments produce the $H = 1$ OHC ground
state excitation. All three are the same phenomenon at
different scales and Hopf charges.
\textbf{Prediction \#162:} plasmoid diameter will scale as
$D_n = D_1 \cdot (j_{1,1}/j_{0,1})^{n-1}$ for Hopf charge
$H = n$, giving a universal quantised diameter spectrum
testable in cavitation experiments.
Zero free parameters.
\smallskip\hrule\smallskip}
\end{center}
\vspace{4pt}
}]

%%----------------------------------------------------------------
\section{Introduction: What Is a Plasmoid?}
%%----------------------------------------------------------------

A plasmoid is a self-contained plasma structure: a doughnut or
toroidal cluster of net protons or net electrons, captured in
a toroidal orbit, self-sustaining through its own
electromagnetic containment field. Plasmoids absorb, store,
and release energy. They are stable against external
perturbations within their containment geometry. They form
naturally in ball lightning, in laboratory arc discharges,
in plasma focus devices, and --- as Malcolm Bendall
demonstrated~\cite{Bendall2022} --- in acoustic cavitation
of water.

The key physical question that plasmoid physics has never
satisfactorily answered is: \textit{why are plasmoids stable?}
What determines their geometry? Why do two counter-rotating
vortices with a zero-point between them form a self-sustaining
structure rather than simply dispersing?

BCT answers this. A plasmoid is not merely a plasma
configuration. It is a macroscopic expression of the
topological structure of the quantum vacuum.

%%----------------------------------------------------------------
\section{The OHC Hopf Charge Framework}
%%----------------------------------------------------------------

\subsection{The Octet-Hopfion Condensate}

The BCT Superfluid Lattice Model identifies the quantum vacuum
ground state as the Octet-Hopfion Condensate (OHC): a
superfluid Planck-scale lattice supporting stable topological
excitations characterised by integer Hopf charge
$H \in \mathbb{Z}$~\cite{BCT_AppJH}.

The Hopf charge is a topological invariant --- it cannot be
changed by continuous deformations of the field. An $H = n$
OHC excitation is topologically protected: it cannot decay
without a topological phase transition. This is why plasmoids
are stable. It is not the electromagnetic containment field
that stabilises them --- it is topology.

\subsection{Hopf charge and plasmoid structure}

The relationship between Hopf charge and plasmoid geometry
is:

\begin{center}
{\small
\begin{tabular}{@{}p{1.2cm}p{1.5cm}p{2.8cm}@{}}
\toprule
$H$ & Vortices & Physical realisation \\
\midrule
0 & None & OHC vacuum ground state \\
1 & 1 & SAFIRE plasma discharge \\
2 & 2 & Standard plasmoid / ball lightning \\
3 & 3 & High-energy ball lightning \\
$n$ & $n$ & General Hopf $n$-plasmoid \\
\bottomrule
\end{tabular}}
\end{center}

\medskip
The $H = 2$ case --- two counter-rotating vortices with a
zero-point at their plane of equilibrium --- is precisely
the structure Bendall describes as the fundamental plasmoid
unit. BCT identifies this as the simplest stable non-trivial
OHC topological excitation.

The ``zero point'' Bendall places at the plane of equilibrium
between the two vortices is, in BCT language, the OHC vacuum
coupling node --- the point where the macroscopic plasmoid
locks into the underlying vacuum topology. This is not
metaphor. It is the physical mechanism of plasmoid stability.

%%----------------------------------------------------------------
\section{Acoustic Cavitation as OHC Excitation}
%%----------------------------------------------------------------

\subsection{The Thunderstorm Generator}

Bendall's Thunderstorm Generator~\cite{Bendall2022} operates
as follows:

\begin{enumerate}\setlength\itemsep{2pt}
  \item Pre-ionised air is introduced into a water-filled
        bubbler chamber.
  \item Microscopic ionised bubbles undergo acoustic
        cavitation --- symmetric collapse generates
        extreme local pressure and temperature.
  \item The collapsing bubble forms a self-contained
        toroidal plasma structure: a plasmoid.
  \item The plasmoid is harvested and fed into a
        conventional engine as an energy source.
\end{enumerate}

BCT provides the mechanism missing from Bendall's own account:
step 3 works because symmetric bubble collapse in an ionised
medium provides exactly the boundary condition for an $H = 2$
OHC topological excitation to nucleate. The two counter-rotating
vortices are the two halves of the collapsing bubble. The zero
point is the OHC vacuum node at the collapse centre.

This is also the mechanism of single-bubble
sonoluminescence~\cite{BCT_L110}: the light flash at the
moment of symmetric collapse is OHC Bessel mode relaxation
radiation. The SBSL spectral peak at 300--400~nm (BCT
Prediction \#154, $<1\%$ error) is the same physical event
as plasmoid nucleation, at lower energy density.

\subsection{The cavitation--plasmoid--sonoluminescence unity}

BCT reveals that three apparently distinct phenomena are the
same OHC event at different energy scales:

\begin{center}
{\footnotesize
\begin{tabular}{@{}lll@{}}
\toprule
Phenomenon & H & Scale \\
\midrule
Sonoluminescence & 2 & low \\
Micro-plasmoid & 2 & medium \\
Ball lightning & 2 & high \\
\bottomrule
\end{tabular}}
\end{center}

\medskip
The geometry is identical. The Hopf charge is identical. Only
the energy scale differs. This is BCT cross-scale universality
applied to plasma physics.

%%----------------------------------------------------------------
\section{Quantised Plasmoid Diameters}
%%----------------------------------------------------------------

The OHC Bessel resonance condition imposes quantisation on
the allowed plasmoid diameters. For Hopf charge $H = n$,
the stable diameter is:
\begin{equation}
  D_n = D_1 \cdot \left(\frac{j_{1,1}}{j_{0,1}}\right)^{n-1}
  = D_1 \cdot 1.5934^{n-1},
\end{equation}
where $D_1$ is the fundamental plasmoid diameter set by the
cavitation bubble size, $j_{1,1} = 3.8317$, and
$j_{0,1} = 2.4048$.

This gives a universal geometric progression:
\begin{align}
  D_1 &= D_0 \quad (H=1) \notag \\
  D_2 &= 1.5934\,D_0 \quad (H=2) \notag \\
  D_3 &= 2.5389\,D_0 \quad (H=3) \notag \\
  D_4 &= 4.0441\,D_0 \quad (H=4). \notag
\end{align}

Ball lightning diameters range from centimetres to tens of
centimetres. If the fundamental scale $D_0 \sim 2$~cm, the
BCT spectrum predicts stable ball lightning at 2, 3.2, 5.1,
and 8.1~cm --- consistent with reported observations of
preferred ball lightning size clusters~\cite{Stenhoff1999}.

%%----------------------------------------------------------------
\section{The SAFIRE Connection}
%%----------------------------------------------------------------

The SAFIRE project (Stellar Atmospheric Function In Regulated
Environments) --- led by Montgomery Childs and Michael Clarage
at Thunderbolts.info --- has produced laboratory plasma
discharge experiments showing stratified glow discharge layers,
anode tufting, and self-organising plasma structures that
replicate solar atmospheric features.

In BCT terms, the SAFIRE discharge produces $H = 1$ OHC
excitations: single-vortex Hopf structures, the simplest
non-vacuum topological state. The stratification layers
correspond to Bessel mode $j_{0,n}$ standing waves in the
discharge column. The anode tufts are $H = 1$ nucleation
sites where the discharge locks into the OHC topology.

The progression SAFIRE ($H=1$) $\to$ Bendall plasmoid (2)
$\to$ ball lightning ($H \geq 2$) $\to$ Valles Marineris
planetary discharge (Letter~117) is a single BCT framework
spanning twelve orders of magnitude in physical scale.

%%----------------------------------------------------------------
\section{Prediction}
%%----------------------------------------------------------------

\noindent\textbf{Prediction \#162} (Quantised plasmoid diameters):
Plasmoid diameters produced by acoustic cavitation at fixed
bubble nucleation size $D_0$ will cluster at the BCT Bessel
spectrum:
\begin{equation}
  D_n = D_0 \cdot 1.5934^{n-1}, \quad n = 1, 2, 3, \ldots
\end{equation}
with clustering width $\sigma < 10\%$ of $D_n$.

\textit{Falsification:} Measure plasmoid diameters in a
controlled cavitation experiment across $>100$ events at
fixed acoustic frequency. If diameter distribution shows
no clustering at the BCT geometric progression: prediction
fails.

%%----------------------------------------------------------------
\section{A Note on Malcolm Bendall}
%%----------------------------------------------------------------

Malcolm Bendall is an Australian inventor and geochemist
who developed his Thunderstorm Generator over more than a
decade~\cite{Bendall2022}. His energy claims --- excess
energy production, over-unity output --- remain unverified
by independent peer review, and BCT takes no position on
those claims.

What BCT does claim is this: the \textit{geometry} of
Bendall's device is correct. The two counter-rotating
vortices, the symmetric bubble collapse, the zero-point
at the plane of equilibrium --- this is $H = 2$ OHC
topology. Bendall arrived at the right geometry through
intuition, observation, and iteration. BCT derives the
same geometry from the structure of the quantum vacuum.

He built a BCT Hopf resonator without knowing why it worked.
That is not a criticism. It is the history of technology.
The Wright brothers did not know why aerofoils generated
lift in terms of the Navier-Stokes equations. They flew
anyway.

%%----------------------------------------------------------------
\begin{thebibliography}{9}

\bibitem{Bendall2022}
M.\ Bendall,
\textit{MSAART Patent Application Notes},
Strike Foundation (2022).
\url{https://www.strikefoundation.earth}

\bibitem{Stenhoff1999}
M.\ Stenhoff,
\textit{Ball Lightning: An Unsolved Problem in Atmospheric
Physics},
Kluwer Academic / Plenum Press (1999).

\bibitem{BCT_AppJH}
M.R.\ Cabri\'{e},
\textit{Appendix JH --- The Hopfion Interior and OHC},
BCT Superfluid Lattice Model (2026).
Zenodo: \href{https://doi.org/10.5281/zenodo.18975018}
{10.5281/zenodo.18975018}.

\bibitem{BCT_L110}
M.R.\ Cabri\'{e},
\textit{Letter 110 --- Vacuum Luminescence},
BCT Superfluid Lattice Model (2026).
Zenodo: \href{https://doi.org/10.5281/zenodo.19177501}
{10.5281/zenodo.19177501}.

\bibitem{BCT_L117}
M.R.\ Cabri\'{e},
\textit{Letter 117 --- The Martian Tetrahedron},
BCT Superfluid Lattice Model (2026).

\bibitem{BCT_GoE}
M.R.\ Cabri\'{e},
\textit{The Geometry of Everything, v4},
BCT Superfluid Lattice Model (2026).
ORCID: 0009-0007-9561-9859.

\end{thebibliography}

\bigskip\noindent
\textit{Acknowledgements.}
Malcolm Bendall --- for building the right geometry before
the theory existed to explain it.
Michael Clarage and Montgomery Childs (SAFIRE) --- for
the laboratory plasma work that BCT now contextualises.
David Talbott and Wal Thornhill --- for the Electric Universe
framework that BCT grounds in quantum vacuum theory.
Michael Steinbacher ($\dagger$~c.2015) --- honoured here.
Jeff, Nic, Nyssa, Tegan, and the Gang Gang Cockatoos.
\textit{The geometry is right. The vacuum insisted.}

\end{document}
